\documentclass[11pt]{article}
\usepackage[a4paper,margin=1in]{geometry}
\usepackage{amsmath,amssymb,amsthm,microtype}
\newtheorem{theorem}{Theorem}[section]
\theoremstyle{remark}\newtheorem{remark}[theorem]{Remark}
\title{Sharp Saturation of the Annular Coefficient-Envelope Rate}
\author{Bin Wang}\date{July 2026}
\begin{document}\maketitle
\begin{abstract}
The mass-limited rate of RH-305 is optimal at the level of its coefficient
information.  We truncate the actual deterministic anchor sequence at order
$N$ and set every later model coefficient to zero.  The resulting sequence
satisfies the same cap envelope with parameter
$M_N=48q^2(q_*/q)^N$, while its annular mismatch is exactly the target tail.
Both its $H^\infty$ and $H^2$ norms are of order
$M_N^{-\kappa(\rho)}/\log M_N$.  Thus neither the power nor logarithmic
factor can be improved from the coefficient envelope alone.  The truncated
sequence is not asserted to be a power-sum sequence of any noisy spectral
complement.
\end{abstract}

\section{Truncated target family}
Fix $0<\rho<\rho_*$ and put $r=q_*/q>1$, $y=q_*\rho<1$.  For $N\ge2$ define
\[
 \tau_n^{(N)}=
 \begin{cases}a_n,&2\le n\le N,\\0,&n>N,\end{cases}
 \qquad
 M_N=48q^2r^N.
\]
The unified target envelope gives, for every $n\ge2$,
\[
 |\tau_n^{(N)}|\le M_Nq^{n-2}.
\]
The corresponding logarithmic mismatch is
\[
 g_N(z)=-\sum_{n>N}\frac{a_n}{n}z^n.
\]

\section{Information-class sharpness}
\begin{theorem}\label{thm:sharp}
For each fixed $0<\rho<\rho_*$ there are constants
$0<c_\rho<C_\rho<\infty$ such that, for
$X=H^\infty(\rho)$ and $H^2(\rho)$,
\[
 c_\rho\frac{y^N}{N}
 \le\|g_N\|_X\le
 C_\rho\frac{y^N}{N}.
\]
Equivalently,
\[
 \|g_N\|_X
 =\Theta_\rho\left(
 \frac{M_N^{-\kappa(\rho)}}{\log M_N}
 \right),
 \qquad
 \kappa(\rho)=\frac{\log(1/y)}{\log r}.
\]
\end{theorem}
\begin{proof}
The upper bounds follow from $|a_n|<48q_*^n$ and the geometric estimates
\[
 \sum_{n>N}\frac{y^n}{n}=O_\rho(y^N/N),\qquad
 \left(\sum_{n>N}\frac{y^{2n}}{n^2}\right)^{1/2}
 =O_\rho(y^N/N).
\]
For the lower bound, let $n_N$ be the first odd integer greater than $N$;
then $n_N\le N+2$ and
\[
 \frac{|a_{n_N}|\rho^{n_N}}{n_N}
 =\frac{y^{n_N}}
 {n_N(1+\lambda^{-n_N})}
 \ge c_\rho\frac{y^N}{N}.
\]
This one coefficient bounds the $H^2$ norm.  For $H^\infty$, take the odd
projection $(g_N(z)-g_N(-z))/2$, whose norm is at most $\|g_N\|_\infty$;
at positive radius all its odd target-tail terms have the same sign, so the
same first term gives the lower bound.  Finally
$M_N=48q^2r^N$ converts $y^N/N$ to the displayed mass scale.
\end{proof}

\section{Protocol and firewall}
The computation records rigorous lower and geometric upper bounds for the
infinite constant-$48$ model tails at $\rho=1.41$ and $N=20,40,80$; it does
not truncate an infinite sum silently.  The numerical model checks the rate
only.

\begin{remark}
The family $\tau^{(N)}$ is a coefficient array satisfying the same envelope.
No multiset in $|\mu|\le q$, no noisy operator, and no physical spectral
power-sum realization is constructed.  Therefore the theorem proves
sharpness only for the coefficient-envelope information class, not for the
actual noisy complement.  Gates A--E remain false/open.
\end{remark}
\end{document}
